% Pais gap loop -- mathematics run sheet, part 8: run-10 items (loaded by \input from part 2, after part 7).
% Items keep their chain numbers (counters set below). Numbers: run10_numbers.json from run10_math.py; canary10.py re-derived 15 run-9 numbers.
%=====================================================================================================
\section*{Run 10 additions to chains I, III and V}
\addcontentsline{toc}{section}{Run 10 additions: I.33--I.35, III.9--III.10, V.13}
\noindent\textbf{Notes on earlier items found in run 10.}
\begin{itemize}
\item \emph{I.32 (emissivity).} The best cryogenic black coatings keep $\varepsilon\ge0.95$ only down to about 30 K \cite{birb}; others fall
steeply below 50--100 K because they turn partly transparent at long wavelengths \cite{acktar}. At 5--25 K the rotor radiates at
120--580 $\mu$m, so $\varepsilon$ there is unmeasured and probably small. I.33 shows this matters little: the allowed ripple scales as $\varepsilon^{1/3}$.
\item \emph{I.30 (sliding debris).} Ricochet data exist only for flat plates: at oblique incidence the largest ricochet fragments keep 24--36\,\% of
the impact speed \cite{schonberg}. Nobody has tested debris sliding along a concave wall at $\ge3$ km/s. I.34 designs that test.
\item \emph{I.32 (thick film).} Long thick-film tapes vary in $I_c$ along their length at 4.2 K, 30 T, and thick films on thin stabilizers show
flux jumps \cite{jaroszynski25}. I.35 sizes the second effect against temperature.
\end{itemize}
Numbers below come from \texttt{run10\_math.py}; the run opened with a blind re-derivation of fifteen run-9 numbers (\texttt{canary10.py}, 15 of 15
within $10^{-4}$).

\setcounter{section}{1}\setcounter{subsection}{32}
\subsection{The field quality B6's stator needs (B6)}
\meta{[P] run 10 $\cdot$ gap 2 (B6) $\cdot$ from I.27, I.32; evidence: Iseult \cite{iseult}, a 9.4 T all-REBCO NMR magnet \cite{rebconmr}}
\begin{proposition}\label{prop:ripple}
Let the rotor see an azimuthal field ripple $\Delta B$ at harmonic $n$ of its rotation, with hysteresis loss $P=nP_0(\Delta B/B_0)^3$
($P_0=1$ W at $B_0=6.3$ mT, I.27), and shed heat only by radiation (emissivity $\varepsilon$, area $A$) to a wall at $T_s$. It stays at or below
the temperature $T_m$ its tape tolerates (I.32) if and only if
\begin{equation}
\Delta B\le B_0\left[\frac{\varepsilon\sigma A\,(T_m^4-T_s^4)}{nP_0}\right]^{1/3}.\label{eq:ripple}
\end{equation}
The bound grows as $\varepsilon^{1/3}$ and falls as $n^{-1/3}$.
\end{proposition}
\begin{proof}
In steady state the absorbed ripple loss equals the radiated power, $nP_0(\Delta B/B_0)^3=\varepsilon\sigma A(T_r^4-T_s^4)$; $T_r$ increases
with $\Delta B$, so $T_r\le T_m$ is the stated inequality.
\end{proof}
\inputs{$A=2.5635$ m$^2$, $B=30$ T, $T^*=22$ K (21--23); $T_m$ from $i\le0.506$: 9.2 K (9.0--9.4) for thick film on a 0.1 mm basis,
24.5 K (23.5--25.4) on a 50 $\mu$m basis; $\varepsilon$ 0.02--0.95; $T_s$ = 4.2, 15, 20 K.}
\nums{At $T_s=4.2$ K and $n=1$: 0.17--0.62 mT (5.7--21 ppm of 30 T) for the 0.1 mm basis and 0.64--2.31 mT (21--77 ppm) for the 50 $\mu$m basis,
across $\varepsilon=0.02$--0.95. A 48-fold range in $\varepsilon$ moves the bound only 3.6$\times$. With $T_s=20$ K (50 $\mu$m basis): 0.52--1.90 mT
(17--63 ppm). Evidence: Iseult (NbTi, 11.72 T) measured 0.88 ppm peak to peak on a 22 cm sphere with passive shims only \cite{iseult};
at 30 T that is 26 $\mu$T, a loss of $7\e{-8}$ W, and the rotor stays at 4.20--4.23 K. A 9.4 T all-REBCO NMR magnet differed from its design
by more than 100 ppm before shimming (screening currents, tape thickness) and reached sub-ppm only after passive and active shimming
\cite{rebconmr}; 100 ppm of 30 T (3 mT) would heat the rotor to 35--62 K and no tape survives. Simulated screening-current fields reach 7.8\,\%
of the centre field in high-field REBCO coils \cite{scfsim}.}
\verdict{The emissivity question folds into a field-quality condition: B6's stator must hold its azimuthal field to about 6--80 ppm along the
rotor's 12.6 m path, depending on the tape. That lies between unshimmed REBCO coils ($\ge100$ ppm) and shimmed NMR/MRI magnets ($\le1$ ppm)
(C: met in small volumes, never along a 4 m ring). Shimming and current overshoot are known remedies; drift of screening currents is the open
item.}
\refs{\cite{iseult,rebconmr,scfsim,birb,acktar}}

\subsection{A tangential fragment test with the craft's own wall pressure (B6)}
\meta{[P] run 10 $\cdot$ gap 2 (B6) $\cdot$ tests I.30; facilities \cite{swrilgg}; no prior test of sliding on a concave wall found}
\begin{proposition}[similarity and the capstan law]\label{prop:capstan}
(a) A coupon cut from the rotor's cross-section, with height $h_c\le h$, the rotor's radial thickness and any length $\ell$, moving at $v$ along
a concave track of radius $R$, presses on the track with the craft's own wall pressure $p=\lambda v^2/(Rh)$, where $\lambda=M/2\pi R$.
(b) If it slides with friction coefficient $\mu$ and loses no mass, its speed after turning through angle $\theta$ is $v_0e^{-\mu\theta}$, so
\begin{equation}
\mu=-\frac{1}{\theta}\ln\frac{v_{\rm out}}{v_{\rm in}},\qquad \Delta E/E=1-e^{-2\mu\theta}.\label{eq:capstan}
\end{equation}
\end{proposition}
\begin{proof}
(a) The coupon's mass is $\lambda(h_c/h)\ell$ and it needs the inward force $\lambda(h_c/h)\ell v^2/R$, spread over the contact area $h_c\ell$:
the pressure is $\lambda v^2/(Rh)$ for any $h_c$ and $\ell$. (b) The track supplies the normal force $N=mv^2/R$, friction $\mu N$ slows the
coupon, $m\,dv/dt=-\mu mv^2/R$; with $ds=v\,dt=R\,d\theta$ this is $dv/d\theta=-\mu v$.
\end{proof}
\inputs{$\lambda=1.78$ kg/m, $h=0.1$ m, $R=2$ m, $v=6.28$ km/s; mean rotor density 8,000 kg/m$^3$ (assumed), so radial thickness 2.2 mm;
coupon $20\times25\times2.2$ mm; sabot 20 g; SwRI's 38 mm two-stage gun: 7 km/s top, 0.5 kg at 3 km/s \cite{swrilgg}.}
\nums{Coupon 8.9 g, normal force 176 kN, wall pressure 351 MPa (the craft's), kinetic energy 176 kJ; with sabot 570 kJ at 6.28 km/s, a quarter
of the gun's 2.25 MJ at 3 km/s. Over a 30$^\circ$ arc (1.05 m, 0.17 ms) $v_{\rm out}/v_{\rm in}$ = 0.949, 0.855, 0.770 for $\mu$ = 0.1, 0.3, 0.5;
1\,\% velocimetry resolves $\mu$ to $\pm0.027$. Over a full turn the same $\mu$ leave 0.53, 0.15, 0.04 of the speed. Entering from a gap of
1, 5, 20 mm reproduces the radial impacts of I.30: 199, 443, 882 m/s. At 3 km/s the coupon carries 4.5 MJ/kg (0.6 of vaporization) and presses
80 MPa; a 0.25 m coupon then weighs 89 g (400 kJ).}
\verdict{A 9 g coupon on a 2 m-radius track reproduces the craft's wall pressure exactly, and its exit speed measures the friction that
sets I.30's slowing time. Two speeds bracket the vaporization threshold (3 and $\sim$6.3 km/s). I.30 fails if the load exceeds $\lambda v^2/R$
by much beyond the entry impact, or if debris leaves the track inward and strikes elsewhere with more than the predicted energy. Facility: C
(energy within reach; launching a flat coupon edge-on at 6 km/s is unproven).}
\refs{\cite{swrilgg,schonberg}}

\subsection{Flux-jump margin against temperature, and a warm design point (B6)}
\meta{[P] run 10 $\cdot$ gap 2 (B6) $\cdot$ criterion from \cite{wilsonmag}; observation \cite{jaroszynski25}; heat capacity model \cite{kittel}}
\inputs{Adiabatic stability of a slab of half-thickness $a$ \cite{wilsonmag}: $\beta=\mu_0J_c^2a^2/(C\,T^*)<3$, with $T^*$ replacing
$T_c-T_0$ for $J_c\propto e^{-T/T^*}$. Thick film: $J_c=36.3$ MA/cm$^2$ (30 T, $B\parallel ab$, 4.2 K), $a=2$ $\mu$m. Heat capacity of the
Hastelloy-dominated tape $C=\gamma T+aT^3$ with $\gamma=0.119$ J kg$^{-1}$K$^{-2}$ (Ni-like) and $\theta_D=400$ K (assumed, $\pm50$\,\%);
REBCO alone 630 J m$^{-3}$K$^{-1}$ at 4.2 K (assumed).}
\nums{$\mu_0J_c^2a^2=6.6\e5$ J/m$^3$. At 4.2 K: $\beta=48$ for the film alone, 0.50 and 0.25 if the heat spreads through a 50 $\mu$m or 0.1 mm tape.
Relative to 4.2 K, $\beta(T)/\beta(4.2)=e^{-2(T-4.2)/T^*}C(4.2)/C(T)$ = 0.19, 0.057, 0.020, 0.007 at 10, 15, 20, 25 K. A warm design point,
wall at 15--20 K and rotor below 24.5 K on the 50 $\mu$m basis, tolerates ripple of 0.71--0.82 mT at $\varepsilon=0.05$ and 1.5--1.8 mT at
$\varepsilon=0.5$ (I.33).}
\verdict{Whether the film flux-jumps at 4.2 K depends on how fast heat leaves the 4 $\mu$m film, which only a measurement settles; flux jumps have
been seen in thick films on thin stabilizers \cite{jaroszynski25}. Running warm cuts the margin parameter 18--50 fold at 15--20 K while the
50 $\mu$m-basis tape still carries the load, and a 20 K wall matches the cryocooler practice of the 20 T SPARC model coil \cite{tfmc}. The
conductor test should therefore ramp the field at spin-up rates at 15--25 K as well as at 4.2 K (D).}
\refs{\cite{wilsonmag,kittel,jaroszynski25,tfmc}}

\setcounter{section}{3}\setcounter{subsection}{8}
\subsection{High-frequency gravitational-wave cavities as a test of EM-to-GW claims}
\meta{[P] run 10 $\cdot$ gap 5 and side thread $\cdot$ from ledger 1 and 7; detectors \cite{hfgwcav}}
\inputs{Plane-wave flux $S=c^3\omega^2h^2/(32\pi G)$; detector reach: $h\approx10^{-21}$ at 300 MHz, $10^{-20}$ at 3 GHz, $3.9\e{-21}$ at 5.311 GHz
\cite{hfgwcav}; Pais's flux formula 1.18$\e{24}$ W/m$^2$ (ledger 1); conversion $4\pi G/(\mu_0c^4)=8.26\e{-38}$ per (T m)$^2$ \cite{gertsenshtein}.}
\nums{The detectors need $S=1.4\e{10}$, $1.4\e{14}$, $6.8\e{13}$ W/m$^2$ respectively. Pais's flux, read as gravitational waves, gives
$h=9.1\e{-15}$, $2.7\e{-15}$, $5.1\e{-16}$ at 0.3, 1, 5.3 GHz: $1.3\e5$ above the 5.3 GHz limit. Known physics for a generous generator
(10 T m, 1 MW) gives $8.3\e{-30}$ W of gravitational waves, $6.6\e{-31}$ W/m$^2$ at 1 m and $h=2\e{-42}$ at 1 GHz, $1.5\e{20}$ below reach.}
\verdict{An existing cavity detector, tuned near a claimed generator's frequency, would see Pais-scale emission at once and known-physics emission
never: a clean null test of any HFGW-generator claim. Stress-energy quadratic in the fields radiates at $2f$, so tuning to $2f$ rejects direct
electromagnetic pickup at $f$; shielding must still exceed the claimed signal ratio.}
\refs{\cite{hfgwcav,gertsenshtein}}

\subsection{Patent sweep 2023--2025: every claim falls under a ledger bound}
\meta{[P] run 10 $\cdot$ side thread $\cdot$ sources \cite{patents25,zhang25pat,oldpatents}}
\inputs{Filings found: US 12,417,870 (gravity field from a varying EM field); US 2025/0055389 (pulsed field to ``decouple the Higgs field''); US
2024/0063731 (inertia attenuator); US 2024/0375795, 2025/0286480, 2025/0296708 (closed electromagnetic drives); US 2025/0132082 (accelerated
electrons radiate); US 12,442,361 (photon drive); older US 3,095,167, 3,610,971, 3,120,363, 3,071,705, 2003/0209637.}
\nums{Each falls in one class. Closed EM systems: thrust $\le P/c=3.34$ nN/W (ledger 17). Radiation or photon drives: exactly $P/c$. Mass or
inertia reduction: field mass is positive and $\le2\e{-9}$ (ledger 16); charge leaves weight unchanged to $2\e{-3}$ (ledger 25). Gravity from
EM fields: $8.26\e{-38}$ per (T m)$^2$ (ledger 7); nothing seen at nN (Tajmar 2024 \cite{tajmar24}). Air-breathing electrostatics: ion wind
(ledger 21). The one checkable number, US 3,095,167 (charge a vehicle against Earth's fair-weather field, $\sim$130 V/m): floating 1 kg needs
0.075 C, but a 1 m sphere at air breakdown (3 MV/m) holds $3.3\e{-4}$ C, 226$\times$ short; it would float 4.4 g.}
\verdict{No new physics thread: the 2023--2025 filings repeat bounded classes, with no data. The side thread's only open edge stays the
high-frequency coupling (K4), testable with balances or a resonant gravity detector (ETH, 42 Hz, 3\,\% agreement with theory \cite{brack22}).}
\refs{\cite{patents25,zhang25pat,oldpatents,tajmar24,brack22}}

\setcounter{section}{5}\setcounter{subsection}{12}
\subsection{FELICE: intracavity fluence against the hold fluence (B7)}
\meta{[P] run 10 $\cdot$ gap 9 (B7) $\cdot$ from V.9, V.12; source \cite{felice}}
\inputs{FELICE first lasing: 10 $\mu$J outcoupled $\approx$ 400 $\mu$J intracavity per micropulse; Rayleigh range 50--82 mm at the user
stations \cite{felice}; $\lambda=30$ $\mu$m (10 THz); hold fluence 0.4 mJ/cm$^2$ \cite{rowe}.}
\nums{Intracavity factor 40. Waist $w_0=\sqrt{\lambda z_R/\pi}$ = 0.69--0.88 mm; peak fluence $2E/\pi w_0^2$ = 33--53 mJ/cm$^2$, 81--133$\times$ the
hold fluence. A film placed 0.58--0.73 m from focus sees the hold fluence over an 8 mm spot.}
\verdict{FELICE's cavity holds far more fluence than B7 needs, so a film could sit well off focus with a large spot and small loss. Its stations
are built for gas-phase samples, with limited access during operation; a cryogenic film station would be new (hardware C).}
\refs{\cite{felice,rowe}}

%=====================================================================================================
\section*{Run 11 notes and the stop rule}
\addcontentsline{toc}{section}{Run 11 notes and the stop rule}
\noindent Run 11 was the last literature pass: kill searches on run 10's conclusions, since-date checks and one bridge (ten searches; numbers
in \texttt{run11\_math.py}; canary: fifteen run-10 numbers re-derived exactly).
\begin{itemize}
\item \emph{I.33, precedent and caution.} Superconducting cyclotrons hold the first azimuthal harmonic at the few-gauss level: 2--3 G already
disturbs the beam at Kolkata; RIKEN's trim coils leave errors below 15 G; a 1 mm trim-rod error at CYCIAE-230 gives 0.24 G \cite{cyclotrons}.
That is the absolute level B6 needs at $n=1$ (0.17--2.3 mT), reached with NbTi and iron at a few tesla. The only REBCO circular-orbit magnet
built for field accuracy, the 0.7 m Ultra-Baby Skeleton Cyclotron (1.6 T design, 540 A, 30 K), quenched at 250 A while charging (46\,\% of its
design current, 36\,\% of $I_c$) and damaged its main coils; screening-current stress and cooldown contraction were blamed \cite{ubsc}. The
technique exists; its transfer to REBCO is unproven and has a failure on record.
\item \emph{I.34, gouging and attitude.} Sled rails gouge above 1.5 km/s and gun barrels above 4 km/s \cite{gouging}; B6's debris at 6.28 km/s
is 1.6--4.2 times faster, so the coupon will gouge its track. The shot must record gouge depth and shape, and the track's material pair should
be chosen from the gouging literature. Entry from a 1, 5 or 20 mm gap meets the wall at 1.8$^\circ$, 4.0$^\circ$ or 8.1$^\circ$; fragment
simulants in a two-stage gun fly with mean pitch and yaw of 4.7--5.5$^\circ$ end-on but 32--43$^\circ$ side-on \cite{fsp}. So the coupon should
fly along its length, with its attitude recorded by flash X-ray.
\item \emph{I.35, buying rules.} Edge-slit thick-film tapes fluctuate in $I_c$ and centre-slit tapes do not; metal layers suppress flux
avalanches in films \cite{jaroszynski25,avalanche}. Buy centre-slit tape with a thicker copper stabilizer.
\item \emph{K4's instrument.} A 42 Hz resonant beam with $Q\approx10^4$ and $\sim$1 pm resolution resolves $F=m\omega^2x/Q\approx$ 7--28 pN for an
effective mass of 1--4 kg (assumed), 36--140 times below a nN balance, at one frequency \cite{brack22,brack23}. A shielded electromagnetic
source whose field energy oscillates at 42 Hz (driven at 21 Hz) is a sharper K4 test than a balance.
\item \emph{X1 and B7.} No outside test and no paper for X1, and no independent driven-K$_3$C$_{60}$ experiment, as of November 2026.
\end{itemize}
\verdict{Run 11 moved no gap: no status, grade or deciding measurement changed, and the tests were refined. Every gap is settled in the loop's
sense: closed (4, 6, 10), against a wall (1, 3, 5, 7), nothing to do (11), or reduced to deciding experiments (2: stator field map,
thick-film ramp, coupon shot; 9: the FEL film test), with gap 8 on a dated re-check (February 2027). The loop stops here. It restarts if a
re-check (gap 8, E-320) finds a new record, an outside group tests Exodus or driven K$_3$C$_{60}$, a REBCO ring is field-mapped, or any
measurement from the test plans comes in.}
\refs{\cite{cyclotrons,ubsc,gouging,fsp,avalanche,brack23}}
